\documentclass[11pt]{article}

\usepackage[margin=1in]{geometry}
\usepackage{amsmath}
\usepackage{amssymb}
\usepackage{amsthm}
\usepackage{enumitem}
\usepackage[colorlinks=true,linkcolor=blue,citecolor=blue,urlcolor=blue]{hyperref}

\newcommand{\F}{\mathbb{F}}
\newcommand{\GF}[1]{\mathbb{F}_{#1}}
\newcommand{\Sym}{S}
\newcommand{\one}{\mathbf{1}}
\newcommand{\Aut}{\operatorname{Aut}}
\newcommand{\Perm}{\operatorname{Perm}}

\title{The $\Sym_5/\GF{5}$ No-Go as an Instance of James's Submodule Theorem:\\
A Trichotomy of Obstruction, Possibility, and Existence\\ for Group-Invariant Codes}
\author{Cheng-Hui Weng\thanks{Literature synthesis assembled from a fan-out, adversarially
verified deep-research pass over the modular representation theory of symmetric groups and
the algebraic-geometry-code secret-sharing literature. The classical core is documented
through graduate lecture notes that quote the named theorems verbatim. Bibliographic
details should be confirmed against the authoritative published sources before formal
citation.}}
\date{30 June 2026}

\begin{document}
\maketitle

\begin{abstract}
The Rocq formalization \texttt{reconstruct/s5\_nogo.v} proves, by an elementary
difference-vector argument, that the natural permutation module $\GF{5}^5$ of the
symmetric group $\Sym_5$ has invariant submodules of dimension only $0,1,4,5$, so no
secret-encoding invariant code of the gap-window dimensions $3$ or $4$ exists. This note
records that the result is not a single witness but the smallest instance of a named
general rule. The obstruction is governed by James's Submodule Theorem for Young
permutation modules. For the natural $\Sym_n$ module the realizable submodule dimensions
are $\{0,1,n-1,n\}$ in every characteristic, since the difference-vector argument is
characteristic-free and the standard module is the irreducible piece of dimension $n-1$, so
the $\Sym_5$ gap window $\{3,4\}$ is infeasible over every field. The divisibility condition
$p \mid n$ controls only whether the structure is uniserial or split, not the available
dimensions, so the obstruction is characteristic-independent and Maschke does not remove
it. Algebraic-geometry-code secret sharing with automorphism-transfer supplies a
constructive existence rule for suitable groups, with the threshold gap governed by the
genus rather than by a submodule lattice, but for the natural $\Sym_5$ action that route is
obstructed as well.
\end{abstract}

\section{The question and the verdict}

A code-based (Massey) secret-sharing scheme that recovers a secret needs a $G$-invariant
submodule $W$, a linear code, of a representation module over a field $F$, of a dimension
forced into a threshold gap window. The $\Sym_5/\GF{5}$ instance acts on $\GF{5}^5$ by the
natural permutation module, with the secret in a fixed coordinate adjoined to the five
shuffled shares. The formalization establishes that the available invariant submodule
dimensions are exactly $\{0,1,4,5\}$, so the intermediate dimensions $2$ and $3$ are
absent and the gap-window code cannot be built.

The verdict of the present survey is that this is a rule, and the rule has a name. The
hand-proof in \texttt{s5\_nogo.v} re-derives the same $\{0,1,4,5\}$ lattice that James's
Submodule Theorem forces. Each of the three tiers below, obstruction, possibility, and
constructive existence, corresponds to a named theorem.

\section{Tier 1: the obstruction rule (which groups cannot)}

The available invariant-code dimensions are precisely the submodule dimensions of the
permutation module, and those are controlled by a single characteristic-free theorem.

\paragraph{James's Submodule Theorem \cite{wildon-symnotes,tomczak-symnotes}.}
For any field $F$ and any partition $\lambda$ of $n$, every $F\Sym_n$-submodule $U$ of the
Young permutation module $M^\lambda$ either contains the Specht module $S^\lambda$ or is
contained in its orthogonal complement $(S^\lambda)^\perp$ under the tabloid bilinear
form. No submodule lies strictly between $S^\lambda$ and $(S^\lambda)^\perp$.

For the natural module $M^{(n-1,1)}$, the relevant object for $\Sym_n$ acting on
$\GF{p}^n$, one has $\dim M = n$, $\dim S^{(n-1,1)} = n-1$ for the sum-zero subspace $P_0$,
and $\dim (S^{(n-1,1)})^\perp = 1$ for the all-ones line.

\paragraph{What the divisibility condition controls.}
The Specht module $S^{(n-1,1)}$ is irreducible if and only if $\operatorname{char} F = 0$
or $\operatorname{char} F \nmid n$ \cite{wildon-symnotes}. This is Peel's theorem on hook
partitions: $S^{(n-1,1)}$ has exactly two composition factors when $p \mid n$ for odd $p$,
and is irreducible otherwise. When $p \mid n$, the all-ones vector has coordinate sum
$n \equiv 0$, hence lies inside the sum-zero space, so
$(S^{(n-1,1)})^\perp \subset S^{(n-1,1)}$ and the lattice collapses to the uniserial chain
\[
  0 \;\subsetneq\; \langle \one \rangle \;\subsetneq\; P_0 = S^{(n-1,1)}
    \;\subsetneq\; M^{(n-1,1)},
  \qquad \dim \in \{0,\,1,\,n-1,\,n\}.
\]
This is the $\{0,1,4,5\}$ chain at $n=5$, and the rigidity that any nonconstant vector
forces the dimension from at most $1$ straight to at least $n-1$. This rigidity is
characteristic-free: the difference-vector argument uses no field hypothesis, so the
submodule dimensions are $\{0,1,n-1,n\}$ over every field. The divisibility $p \mid n$
governs only whether $\langle \one\rangle$ lies inside $P_0$, that is uniserial versus
split, not which dimensions occur. The simple modules are
fixed as well: in prime characteristic $p$ the simple $F\Sym_n$-modules are
$D^\lambda = S^\lambda / (S^\lambda \cap (S^\lambda)^\perp)$ for $p$-regular $\lambda$,
James's 1976 classification recorded in \cite{wildon-symnotes}.

\paragraph{Naming the $\Sym_5/\GF{5}$ module precisely.}
The module is the natural permutation module $M^{(n-1,1)}$, and the precise modular
phenomenon is that when $\operatorname{char} F \mid n$ the $(n-1,1)$-Specht module is a
submodule but not a direct summand of $M^{(n-1,1)}$ \cite{tomczak-symnotes}. Maschke's
complement fails exactly there, and the intermediate piece does not split off.

\paragraph{Generalization.}
The no-go holds for $\Sym_n$ on $\GF{p}^n$ exactly when $p \mid n$, and then the entire
window of intermediate dimensions $2, 3, \dots, n-2$ is forbidden, not merely a single
gap. The case $n = p = 5$ is the smallest instance. The transposition-difference mechanism
of \texttt{s5\_nogo.v}, in which $2$-transitivity breeds every difference vector
$e_i - e_j$ and spans the dimension-$(n-1)$ sum-zero space, is the concrete shadow of this
theorem.

\section{Tier 2: the possibility rule (which groups can)}

\paragraph{Maschke's theorem \cite{tomczak-symnotes}.}
If $\operatorname{char} F \nmid |G|$, every submodule has a $G$-invariant complement, the
module is semisimple, and every partial sum of irreducible-constituent dimensions is
realized as a submodule. This is the general possibility rule. Whether it returns the
gap-window dimensions depends on the irreducible-constituent dimensions of the particular
module.

\paragraph{The possibility rule does not rescue $\Sym_5$ (correction).}
For the natural module $M^{(n-1,1)}$ the constituents are the trivial module of dimension
$1$ and the standard module of dimension $n-1$, so the realizable submodule dimensions are
$\{0,1,n-1,n\}$ in every characteristic. Maschke changes the structure from uniserial to
split, but it does not change this dimension profile. The symmetric group $\Sym_5$ has no
$3$-dimensional ordinary irreducible at all, its irreducible dimensions being
$\{1,1,4,4,5,5,6\}$, and its standard module is the irreducible of dimension $4$. The
dimension-$3$ object $D^{(4,1)}$, the heart, exists only in characteristic $5$, and even
there only as the subquotient $P_0/\langle \one\rangle$, never as a submodule. The
difference-vector argument of \texttt{s5\_nogo.v} is itself characteristic-free, so the
$\{3,4\}$ gap-window infeasibility for the natural $\Sym_5$ action holds over every field.
An earlier reading, following a comment in the Rocq source, claimed that coprime
characteristic restores dimensions $2$ and $3$. That is wrong for $\Sym_5$ and is corrected
here. The contrast is $\Sym_4$, whose standard module is the irreducible of dimension $3$,
so a dimension-$3$ invariant code does exist there, which is exactly why the no-go needs
five shuffled points rather than four.

\paragraph{Refinement away from the corner \cite{wildon2009mfree}.}
For $n \geq 66$ and $p > 3$, the multiplicity-free permutation characters of $\Sym_n$ are
completely classified, and each summand of the modular permutation module is self-dual and
carries a Specht filtration whose factor dimensions are read off the ordinary character.
The Specht module $S^\lambda$ appears in the filtration if and only if the ordinary
character $\chi^\lambda$ is a constituent of the permutation character. Outside small or
defining characteristic, the modular permutation module retains a clean and predictable
composition series.

\paragraph{Small characteristic genuinely breaks order \cite{wildon2009mfree}.}
The hypotheses $p > 3$ and large $n$ are necessary, not artifacts. In characteristic $3$
the module $F[\Sym_9/P\Gamma L(2,8)]$ has no Specht filtration at all, and in
characteristic $2$ the module $\F_2[\Sym_6/(\Sym_3 \wr \Sym_2)]$ has a Specht filtration
but not one with the factors predicted by its ordinary character. The $\Sym_5/\GF{5}$
uniserial collapse is the cleanest member of this small-and-defining-characteristic family.

\paragraph{The semisimple multiplicity data \cite{jamespeel1979,james1978}.}
The composition-factor multiplicities underlying the realizable submodule dimensions in the
Maschke regime are Littlewood-Richardson coefficients: over a field where the skew Specht
module is semisimple, $V^{\lambda/\mu} \cong \bigoplus_\nu (V^\nu)^{c^\lambda_{\mu\nu}}$.
The James-Peel contribution is the corresponding Specht filtration with the same
multiplicities valid over an arbitrary field, the direct-sum form holding only in the
semisimple regime.

\section{Tier 3: the constructive rule (which groups certainly do)}

The constructive engine abandons the permutation-module lattice and switches to
genus-controlled constructions, in which the threshold gap is a different object.

\paragraph{The AG-LSSS gap is the genus \cite{chen2006aglsss}.}
From a residue code on a genus-$g$ curve, subsets of size below $n-m$ are unqualified and
subsets of size at least $n-m+2g$ are qualified, so the privacy and reconstruction
thresholds are separated by a gap of width exactly $2g$, with $m = \deg G$. The thresholds
are bracketed by the minimum and dual distances,
$n-m \leq d_{\min} \leq d_{\mathrm{cheat}} \leq n-m+2g$. Genus $0$ recovers Reed-Solomon and
hence Shamir, with gap $0$. The gap is driven by the genus, not by a submodule lattice.

\paragraph{Existence over any fixed field \cite{cccx2009}.}
For every fixed finite field there is an infinite family of ideal linear secret-sharing
schemes that is asymptotically good with strong multiplication, built on algebraic-geometry
codes over good towers of function fields. This is an unconditional possibility result, and
it imposes no group symmetry on its own.

\paragraph{The symmetry-transfer rule that guarantees invariance \cite{joynerksir2008}.}
If a divisor $D$ is very ample and the evaluation point set $E$ is large enough, namely
$n > r \cdot \deg D$ for some $r \geq 1$, then the curve-automorphism group fixing $D$ and
$E$ is realized exactly as the permutation automorphism group of the code, and the two
groups are isomorphic, $\Aut_{D,E}(X) \cong \Perm(C)$. Concretely, for a smooth projective
curve of genus $g \geq 2$ with $\deg D \geq 2g+1$ and $|E| \geq (1+g)\deg D$, the code has
$n = \deg E$, $k = \deg D + 1 - g$, $d \geq \deg E - \deg D$, and its permutation
automorphism group is the prescribed curve symmetry. Choosing a curve with the desired
automorphism group, such as the Hermitian, Suzuki, Ree, or modular curves with their large
automorphism groups, transfers that symmetry to the code.

\paragraph{Group-algebra codes \cite{pandeykrishnan2026,garciaclaro2020}.}
The Berman codes and their duals, derived from the group algebra $\F_2[\mathbb{Z}_n^m]$,
have automorphism group the wreath product $\Sym_n \wr \Sym_m$ in product action, with the
related abelian codes realizing the same group for odd $n \geq 5$. A sharp
modular-versus-semisimple existence dichotomy parallels the no-go: assuming the MDS
conjecture, which is a theorem for prime fields, the only nontrivial MDS group codes in a
non-semisimple group algebra $\F_p G$ occur for $G = C_p$ with $p$ odd, where they coincide
with extended Reed-Solomon codes, and outside this cyclic corner an MDS group code forces
$\F_p G$ to be semisimple.

\section{Location of the $\Sym_5/\GF{5}$ case}

The case sits in the multiplicity-free corner of James's Submodule Theorem, with
$n = p = 5$. For the natural $\Sym_n$ action the realizable submodule dimensions are
$\{0,1,n-1,n\}$ in every characteristic, because the difference-vector argument is
characteristic-free and the standard module is the irreducible piece of dimension $n-1$.
The divisibility condition $p \mid n$ controls only the module structure, uniserial when
$p \mid n$ and split when $p \nmid n$. It does not change the available dimensions. So the
$\Sym_5$ gap window $\{3,4\}$ is infeasible over every field, Maschke does not restore it,
and the no-go is characteristic-independent for the natural action. The same holds for
$\Sym_n$ on $\GF{p}^n$ for every $n$, the forbidden interior being the dimensions
$2, \dots, n-2$.

\section{Caveats and refuted framings}

Three points constrain how the result should be stated.

\begin{itemize}[leftmargin=1.4em]
\item The no-go should not be framed as a distributive-lattice, join-irreducible-poset, or
  Ext-quiver result. Several such framings were tested and refuted. The supported mechanism
  is specifically James's Submodule Theorem together with the $(n-1,1)$ irreducibility
  dichotomy, and nothing more exotic.
\item The two notions of gap are different objects. The algebraic-geometry threshold gap is
  the genus window of width $2g$, a curve and code parameter. The representation-theoretic
  gap is the set of missing submodule dimensions. The constructive results show how to
  build symmetric codes with controlled parameters. They do not show the submodule
  obstruction reappearing inside algebraic-geometry codes. The two are independent
  constraints that a recoverable invariant scheme must both satisfy.
\item Source quality varies. The James, Maschke, and Peel core is textbook-standard and
  documented through lecture notes that quote the theorems verbatim, hence high confidence.
  The Berman automorphism result is a single recent preprint, hence medium confidence. The
  MDS group-code dichotomy is conditional on the MDS conjecture, discharged for prime
  fields.
\item For the specific natural-action $\Sym_5$ protocol of \texttt{pgg-smc}, the Tier-3
  constructive routes were evaluated and rejected. Genus $0$ is excluded because
  $|\Sym_5| = 120$ exceeds the Klein finite-subgroup ceiling $60$, Wiman's classification
  rules out genus $3$, and Bring's curve at genus $4$ is only axiomatized while the
  algebraic-geometry recovery route is excluded by the no-go itself. Recovery survives only
  for the small cyclic instances and the intransitive product. A companion note records the
  rejected Bring-curve and genus-$0$ PGL routes with commit and file pointers.
\end{itemize}

\section{Open problems}

\begin{enumerate}[leftmargin=1.6em]
\item A clean criterion for which other $2$-transitive groups, such as projective special
  linear, affine, or sporadic $2$-transitive groups, carry the defining-characteristic
  obstruction on their natural modules. The $2$-transitivity is what breeds the difference
  vectors, but no consulted source pins down the general cross-versus-defining split for an
  arbitrary $2$-transitive group.
\item Whether a nontrivial group symmetry can be imposed on the asymptotically good ideal
  scheme family while preserving asymptotic goodness and strong multiplication. The
  automorphism-transfer theorem supplies the mechanism and the asymptotically good family
  exists, but no consulted source combines them into an asymptotically good symmetric
  algebraic-geometry scheme.
\item A closed form for the forbidden intermediate dimensions of a general Young
  permutation module $M^\lambda$, not merely $M^{(n-1,1)}$, in defining characteristic,
  expressed through the $p$-regular Specht-filtration data.
\item Whether the genus gap of width $2g$ and the submodule-lattice gap admit a unified
  representation-theoretic statement, or are genuinely independent constraints that must
  both be satisfied for a recoverable group-invariant covering scheme.
\end{enumerate}

\bibliographystyle{plain}
\bibliography{20260630-170155-s5nogo-james-submodule-report}

\end{document}
